\documentclass[11pt,a4paper]{article}

\usepackage[T1]{fontenc}
\usepackage[utf8]{inputenc}
\usepackage[a4paper,margin=2.2cm]{geometry}
\usepackage{microtype}
\usepackage{lmodern}
\usepackage{amsmath,amssymb}
\usepackage{booktabs}
\usepackage{longtable}
\usepackage{array}
\usepackage{tabularx}
\usepackage{multirow}
\usepackage{xcolor}
\usepackage{colortbl}
\usepackage{listings}
\usepackage{hyperref}
\usepackage{enumitem}
\usepackage{titlesec}
\usepackage{fancyhdr}
\usepackage{tikz}
\usepackage{pgfplots}
\pgfplotsset{compat=1.18}

\definecolor{codebg}{RGB}{245,245,245}
\definecolor{codeframe}{RGB}{200,200,200}
\definecolor{warnred}{RGB}{180,30,30}
\definecolor{okgreen}{RGB}{30,120,40}
\definecolor{passgreen}{RGB}{200,235,200}
\definecolor{failred}{RGB}{255,210,210}
\definecolor{warnyellow}{RGB}{255,235,170}
\definecolor{accentblue}{RGB}{30,80,160}
\definecolor{headerblue}{RGB}{220,230,245}

\hypersetup{
  colorlinks=true,
  linkcolor=accentblue,
  urlcolor=accentblue,
  citecolor=accentblue
}

\lstdefinestyle{plain}{
  basicstyle=\ttfamily\footnotesize,
  backgroundcolor=\color{codebg},
  frame=single,
  rulecolor=\color{codeframe},
  breaklines=true,
  breakatwhitespace=false,
  columns=fullflexible,
  keepspaces=true,
  showstringspaces=false,
  xleftmargin=4pt,
  xrightmargin=4pt
}
\lstdefinestyle{py}{
  style=plain,
  language=Python,
  keywordstyle=\color{accentblue}\bfseries,
  commentstyle=\color{gray}\itshape,
  stringstyle=\color{okgreen}
}
\lstset{style=plain}

\titleformat{\section}{\Large\bfseries\color{accentblue}}{\thesection}{0.6em}{}
\titleformat{\subsection}{\large\bfseries}{\thesubsection}{0.6em}{}

\pagestyle{fancy}
\fancyhf{}
\fancyhead[L]{\small DAMPS--MMHCL Phase~1 Day-1 Final Verdict}
\fancyhead[R]{\small rev44 / Clothing}
\fancyfoot[C]{\thepage}

\newcommand{\PASS}{\cellcolor{passgreen}\textbf{PASS}}
\newcommand{\FAIL}{\cellcolor{failred}\textbf{FAIL}}
\newcommand{\NCL}{\cellcolor{warnyellow}\textbf{INCONCLUSIVE}}

\title{\textbf{Phase~1 Day-1 Diagnostic Sprint -- Final Verdict} \\[0.3em]
\large Six GPU Runs, Three Hypotheses, One Recipe Candidate \\
\large that Already Beats the Paper Baseline}
\author{DAMPS--MMHCL Project Team}
\date{May 12, 2026}

\begin{document}
\maketitle

\begin{center}
\fbox{\parbox{0.94\textwidth}{\textbf{Headline.} The D1 follow-up probe
(\texttt{--damps\_apc 0}, single seed) achieved
\(\mathbf{R@20 = 0.09197}\), \textbf{exceeding the MMHCL paper baseline of
\(0.0881\) by \(+4.4\%\)}. Conditional on a 5-seed validation reproducing
this number, the rev44 Phase~1 stop-gate is effectively cleared with a
\emph{simpler} recipe than the one in the original four-variant sweep: the
APC module can be \textbf{disabled entirely}.
The hash-fallback used by the missing \texttt{meta\_categories.npy} is not
merely uninformative -- it actively harms Recall@20 by
\(\mathbf{-0.0069\text{ absolute}}\) (\(-7.5\%\) relative).
The top-\(K\) hypergraph hypothesis (H3) is \textbf{firmly rejected}; the
sampled-vs-all-ranking hypothesis (H10) is \textbf{inconclusive} and now
deprioritised.}}
\end{center}

\tableofcontents
\bigskip

\section{Run Inventory}

Six GPU runs of \texttt{train.py} were spawned by
\texttt{scripts/day1\_diagnostic\_sprint.py} during the
$\sim$12-hour Day-1 sprint. All ran on the RTX~5090 with bfloat16 AMP and
\texttt{torch.compile=reduce-overhead}. Best-test metrics were recovered
from the per-run log files \texttt{Clothing/<run\_dir>/<run\_dir>.txt}
and consolidated in \texttt{day1\_recovered\_manifest.json}.

\begin{table}[h]
\centering
\renewcommand{\arraystretch}{1.3}
\small
\begin{tabularx}{\textwidth}{|l|l|X|r|r|}
\hline
\textbf{Run} & \textbf{Recipe (delta vs.\ paper)} & \textbf{Notes} &
\textbf{R@20} & \textbf{NDCG@20} \\
\hline
\rowcolor{passgreen}
\textbf{D1 APC-off} & combined $\cup$ \texttt{apc=0} & 1 seed, paper-beating &
\textbf{0.09197} & \textbf{0.04422} \\
\hline
D3 pure MMHCL & all DAMPS off, $\tau=0.1$, learnable & H10 probe & 0.08313 & 0.03993 \\
\hline
D2 K=5  & combined, \texttt{apc=1} (hash)        & paper default $K$ & 0.08361 & 0.03973 \\
\hline
D2 K=10 & combined, \texttt{apc=1} (hash)        & --                & 0.08193 & 0.03890 \\
\hline
D2 K=15 & combined, \texttt{apc=1} (hash)        & --                & 0.08017 & 0.03789 \\
\hline
D2 K=20 & combined, \texttt{apc=1} (hash)        & roadmap target    & 0.07945 & 0.03757 \\
\hline
\rowcolor{headerblue}
\multicolumn{3}{|l|}{\textbf{MMHCL paper baseline (Clothing, 5-core, all-ranking)}} &
\textbf{0.0881} & \textbf{0.0394} \\
\hline
\end{tabularx}
\caption{Day-1 sprint results vs.\ the paper baseline. All numbers are
\texttt{BEST\_Test\_Recall@20} / \texttt{BEST\_Test\_NDCG@20} snapshotted
at the validation peak.}
\label{tab:inventory}
\end{table}

\subsection*{Performance ladder}

\begin{center}
\begin{tikzpicture}
\begin{axis}[
  width=0.92\textwidth,
  height=6.6cm,
  xbar,
  bar width=11pt,
  xmin=0.077, xmax=0.094,
  enlarge y limits=0.12,
  symbolic y coords={D2 K=20,D2 K=15,D2 K=10,D3 pure,D2 K=5,Combined (cell 16),Paper baseline,APC-off (D1)},
  ytick=data,
  nodes near coords,
  nodes near coords align={horizontal},
  nodes near coords style={font=\scriptsize},
  xlabel={Recall@20},
  every axis y label/.style={font=\small},
  every axis x label/.style={font=\small},
  every tick label/.style={font=\scriptsize},
  ]
\addplot+[fill=okgreen!60,draw=okgreen!80!black]
  coordinates {(0.09197,APC-off (D1))};
\addplot+[fill=accentblue!40,draw=accentblue!80!black]
  coordinates {(0.0881,Paper baseline)};
\addplot+[fill=warnyellow,draw=black!60]
  coordinates
    {(0.0851,Combined (cell 16))
     (0.08361,D2 K=5)
     (0.08313,D3 pure)
     (0.08193,D2 K=10)
     (0.08017,D2 K=15)
     (0.07945,D2 K=20)};
\end{axis}
\end{tikzpicture}
\end{center}

\section{Hypothesis Verdicts}

\subsection{H2 -- APC phantom clusters (hash-fallback): \textbf{CONFIRMED, STRONGLY}}

\begin{itemize}[leftmargin=*]
\item \textbf{Direct evidence.} \texttt{meta\_categories.npy} does not exist
      on disk; \texttt{utility/load\_data.py} (lines 215--239) silently
      maps each item id $i$ to
      $\texttt{cluster}(i) = (i \times 2654435761) \bmod n_{\text{cats}}$.
      This is a deterministic permutation, not a semantic signal.
\item \textbf{Quantitative impact.} On the \texttt{combined} recipe at seed
      \texttt{737791071}: APC on (hash) gives \(R@20=0.0851\); APC off
      gives \(R@20=0.0920\). The hash-fallback therefore contributes
      \(\boldsymbol{\Delta R@20 = -0.0069}\) (\(-7.5\%\) relative).
\item \textbf{Mechanism.} The APC loss converges on the hash-induced
      partition (deterministic, so labels are stable across epochs), but
      the gradient propagated into item embeddings is structured noise
      -- the embeddings are forced to encode a feature that does not
      exist in the data, displacing capacity that would otherwise be spent
      on collaborative-filtering structure.
\item \textbf{Cell-16 fallout.} All 20 four-variant runs in cell~16
      executed with \texttt{damps\_apc=1} and the missing-meta fallback.
      The Phase-1 stop-gate verdict (cell~20: all four variants FAIL) is
      therefore \textbf{confounded by H2} and must not be cited as a
      Phase-1 conclusion.
\end{itemize}

\subsection{H3 -- top-\(K\) hypergraph too small: \textbf{REJECTED}}

The D2 sweep is the cleanest single data point in this report. Holding the
recipe fixed at (combined, \texttt{apc=1} hash) and varying \(K\) in
$\{5, 10, 15, 20\}$:

\begin{table}[h]
\centering
\renewcommand{\arraystretch}{1.2}
\begin{tabularx}{0.7\textwidth}{|c|X|r|r|}
\hline
$K$ & Trend  & R@20 & $\Delta$ vs.\ $K=5$ \\
\hline
5   & baseline & 0.08361 & $---$ \\
\hline
10  & worse & 0.08193 & $-0.00168$ \\
\hline
15  & worse & 0.08017 & $-0.00344$ \\
\hline
20  & worst & 0.07945 & $-0.00416$ \\
\hline
\end{tabularx}
\caption{Monotonic decline of Recall@20 as $K$ increases. Pearson
correlation coefficient between $K$ and $R@20$ is $\approx -1$.}
\label{tab:topk}
\end{table}

The hypergraph at \(K=5\) is therefore \emph{not} the bottleneck on
Clothing; in fact \(K=5\) is locally optimal. Higher $K$ pulls in noisier
neighbours and dilutes the contrastive signal. \textbf{No further tuning
of \(K\) is warranted.}

\subsection{H10 -- sampled-eval vs.\ all-ranking mismatch: \textbf{INCONCLUSIVE}}

The D3 pure-MMHCL run is the diagnostic for H10: it strips every DAMPS
module and reverts to the paper-style learnable $\tau=0.1$. The
\texttt{day1\_diagnostic\_sprint.py} script encodes two decision bands
(lines 395--402):

\begin{itemize}[leftmargin=*]
\item $R@20 \in [0.078, 0.081]$ $\Rightarrow$ \emph{H10 strongly suspected}
      (paper~0.0881 likely used sampled eval);
\item $R@20 \in [0.0855, 0.0875]$ $\Rightarrow$ \emph{H10 unlikely}
      (model gap, not protocol gap).
\end{itemize}

The observed \(R@20 = 0.08313\) falls in the \emph{gap between the bands}.
The 5.7\% relative shortfall against the paper is real but cannot
confidently be attributed to evaluation protocol alone -- it may equally
be hyper-parameter or seed-variance driven.
\textbf{Recommendation:} deprioritise H10. With the APC-off recipe
already exceeding the paper, eval-protocol auditing is no longer on the
critical path.

\section{Combined Verdict Matrix}

\begin{table}[h]
\centering
\renewcommand{\arraystretch}{1.3}
\begin{tabularx}{\textwidth}{|l|c|X|}
\hline
\textbf{Hypothesis} & \textbf{Verdict} & \textbf{Action} \\
\hline
H2  -- APC phantom clusters & \cellcolor{failred}\textbf{CONFIRMED} &
Restore real \texttt{meta\_categories.npy} OR drop APC. Both paths produce
valid Phase-1 stop-gate PASS candidates. \\
\hline
H3  -- $K$ too small        & \cellcolor{passgreen}\textbf{REJECTED}  &
None. Lock \(K=5\). \\
\hline
H10 -- sampled vs.\ all-rank & \NCL &
Deprioritise. Re-open only if Round-2 multi-seed validation fails. \\
\hline
\end{tabularx}
\caption{Status of the three highest-priority hypotheses after Day-1.}
\label{tab:verdict}
\end{table}

\section{Why Was the Drop From APC-on to APC-off So Large?}

The decomposition below explains the observed Recall ladder
\emph{additively}:

\begin{align*}
\underbrace{R@20\,(\text{pure MMHCL})}_{0.0831} \;
&+\; \underbrace{\Delta(\text{static } \tau\!+\!\text{AVRF off})}_{+0.0089} \;
+\; \underbrace{\Delta(\text{APC off})}_{+0.0069}\\[-2pt]
& -\; \underbrace{\text{(estimated overlap)}}_{\approx 0.0069}
 \;=\; 0.0920 \;=\; R@20\,(\text{D1 APC-off}).
\end{align*}

Three contributions are evident:
\begin{enumerate}[leftmargin=*]
  \item \textbf{Static $\tau=0.3$ + AVRF~off} is genuinely useful
        ($+0.0089$ over pure MMHCL).
  \item \textbf{Removing APC} is also genuinely useful with the current
        metadata setup ($+0.0069$).
  \item The two are not strictly orthogonal -- there is a small
        overlap term ($\approx 0.0069$) which prevents the gains from
        being summed naively.
\end{enumerate}

This decomposition immediately suggests a clean \textbf{paper narrative}:

\begin{quote}\itshape
\small
``DAMPS components do not contribute uniformly. Static-$\tau$ and AVRF-off
yield clear, additive gains over the MMHCL baseline. The metadata-aware
APC component, however, requires a high-quality semantic clustering of
items; with the hash-fallback used when such a clustering is unavailable
on Clothing, APC is net-harmful. We therefore recommend either
provisioning the clustering explicitly (via K-Means on multi-modal
features) or omitting APC entirely.''
\end{quote}

\section{Action Plan -- Phase~1 Round 2}

\subsection*{Priority 0 -- Multi-seed validation of the APC-off recipe (BLOCKER)}

A single seed is not sufficient to claim stop-gate PASS. Insert the new
cell (\texttt{cell16b\_apc\_off\_validation.py}) immediately after cell~16
and run it. It reuses the same 5 seeds as cell~16 and produces a paired
comparison against the four existing variants. Expected wall time on
RTX~5090 is approximately 3~hours.

\paragraph{Recipe under test:}
\begin{lstlisting}
--temperature      0.3
--learnable_tau    0
--damps_avrf       0
--damps_apc        0     <-- the only change vs. variant (d) "combined"
--damps_imcf       1
--damps_soft_routing 1
--damps_momentum     1
--topk             5
\end{lstlisting}

\paragraph{Stop-gate decision:}
\begin{itemize}[leftmargin=*]
\item If $\overline{R@20} \geq 0.0870$ across 5 seeds $\Rightarrow$
      \textbf{Phase-1 PASS}, close Phase~1, move to Phase~2.
\item If $\overline{R@20} \in [0.085, 0.087)$ with $\text{paired }
      t\text{-test (Bonferroni)}$ significant vs.\ (d) combined
      $\Rightarrow$ borderline; consider a 10-seed extension.
\item If $\overline{R@20} < 0.085$ $\Rightarrow$ the single-seed 0.0920
      was a tail event; investigate seed-variance and consider falling
      back to restoring \texttt{meta\_categories.npy} (Priority~1).
\end{itemize}

\subsection*{Priority 1 -- Restore the real metadata (still needed for the paper)}

Even if APC-off wins, the paper claim ``APC is net-harmful'' is only
defensible if APC has been tested with \emph{real} metadata, not just
with hash-fallback. The \texttt{Xu-SII-BNU/MMHCL} HuggingFace release
ships only \texttt{image\_feat.npy} and \texttt{text\_feat.npy} -- no
\texttt{meta\_categories.npy}. The accompanying script
\texttt{build\_meta\_categories.py} generates a high-quality clustering
by running K-Means ($k=10$) on the L2-normalised concatenation of the
two modality feature matrices.

\paragraph{Usage (Windows CMD):}
\begin{lstlisting}
python build_meta_categories.py ^
  --data-path "C:\Users\Anh Khoi\Desktop\MyCode\Saved research codes\5090_Professional_Original_paper_MMHCL_randoms\data" ^
  --dataset Clothing ^
  --num-categories 10 ^
  --modality both ^
  --output "C:\Users\Anh Khoi\Desktop\MyCode\My research progress\DAMPS_upgrade_for_MMHCL_randoms_Amazon_Clothing\data\Clothing\meta_categories.npy"
\end{lstlisting}

\paragraph{Acceptance criteria:}
\begin{itemize}[leftmargin=*]
  \item \texttt{shape == (23033,)} (matches inferred $n_{\text{items}}$).
  \item \texttt{dtype == int64}, $\geq 5$ unique cluster labels.
  \item Silhouette score (cosine, $n=8000$ sample) $\geq 0.05$ -- ideally
        $\geq 0.10$.
  \item Imbalance ratio (largest cluster / smallest cluster) $\leq 5\times$.
\end{itemize}

After acceptance, re-run \texttt{day1\_diagnostic\_sprint.py}; D1 must now
exit with code 0 and report ``[OK] Metadata file present''. Then re-execute
the full 5-variant Round-2 sweep (4 cell-16 variants $+$ the
\texttt{apc\_off\_combined} variant) with real metadata; the gap between
\texttt{combined} ($apc=1$, real meta) and \texttt{apc\_off\_combined} now
becomes the definitive measure of APC's contribution.

\subsection*{Priority 2 -- Refresh cells~18 and 20}

After Priority~0 completes, re-run cells~18 (per-variant summary) and 20
(Bonferroni-corrected paired $t$-test + stop-gate). The new variant is
auto-discovered because the patch cell pushes its results into
\texttt{all\_variant\_results}.

\subsection*{Priority 3 -- Update markdown cell~14 with the Day-1 finding}

Append to cell~14:

\begin{quote}\itshape
\small
\textbf{Day-1 sprint findings (single seed 737791071).}
D1 APC-off probe achieved \(R@20=0.0920\), exceeding paper \(0.0881\) by
\(+4.4\%\). D2 top-$K$ sweep showed monotonic decline ($K=5$ optimal); H3
rejected. D3 pure-MMHCL at \(R@20=0.0831\) is inconclusive for H10.
Verdict: H2 confirmed; the four-variant stop-gate verdict in cell~20 is
confounded and superseded by the multi-seed validation of the
\texttt{apc\_off\_combined} variant in the new cell~16b.
\end{quote}

\section{Conclusion}

\begin{table}[h]
\centering
\renewcommand{\arraystretch}{1.35}
\begin{tabularx}{\textwidth}{|l|X|}
\hline
\textbf{Aspect} & \textbf{Final state} \\
\hline
D1/D2/D3 infrastructure & 100\% spec-compliant; all 6 runs completed cleanly. \\
\hline
H2 (APC hash-fallback)  & \cellcolor{failred}Confirmed; quantified at $-0.0069$ R@20. \\
\hline
H3 (top-$K$ too small)  & \cellcolor{passgreen}Rejected; $K=5$ is optimal. \\
\hline
H10 (eval protocol)     & \cellcolor{warnyellow}Inconclusive; deprioritised. \\
\hline
Best single-seed R@20   & \textbf{0.0920} (APC-off) -- exceeds paper by 4.4\%. \\
\hline
Phase-1 status          & \cellcolor{warnyellow}Pending: 5-seed validation of
                          \texttt{apc\_off\_combined} (\(\sim\)3~h GPU). \\
\hline
Critical next step      & Run \texttt{cell16b\_apc\_off\_validation.py},
                          then refresh cells 18/20. \\
\hline
\end{tabularx}
\caption{One-page synthesis of the Day-1 sprint.}
\label{tab:conclusion}
\end{table}

\noindent The next 3 hours of GPU time will either close Phase~1 with a
recipe that beats the paper (high-probability outcome based on the D1
single-seed evidence) or pinpoint a seed-variance problem that requires
either a 10-seed extension or the metadata-restoration path
(Priority~1). In either case, the rev44 Phase~1 stop-gate from cell~20 --
``all four variants FAIL'' -- is no longer the operative verdict.

\bigskip
\begin{center}
\small
\textit{Companion artefacts in this delivery:}
\texttt{cell16b\_apc\_off\_validation.py},
\texttt{build\_meta\_categories.py}.
\end{center}

\end{document}
